\section{Affinoid spaces}

    Fix a non-archimedean completed valued field \((\kk,|\cdot|)\) throughout this section.
    Let \(A\) be a \(\kk\)-affinoid algebra, and let \(X = \scrM(A)\) be the associated affinoid space.
    It should be viewed as the non-archimedean analogue of an affine variety in algebraic geometry or an affine scheme in scheme theory.
    The goal of this section is to define a suitable Grothendieck topology on \(X\) such that we can define a structure sheaf of analytic functions on \(X\), 
    and then study its basic properties.



\subsection{Affinoid domains}

    On a general open subset \(U \subseteq X\), the ring of analytic functions on \(U\) may behave badly.

    \Yang{To be continued.}
    % Consider \(X = \scrM(A)\) with \(A = \kk\{T_1,\ldots,T_n\}\).
    % \Yang{Not every open subset of \(X\) gives an affinoid space, that is, the completion of the ring of analytic functions on that open subset is not necessarily an affinoid algebra.}
    % \Yang{Right? example?}

    \begin{definition}\label{def:affinoid_domains}
        A closed subset \(V \subseteq X\) is called an \emph{affinoid domain} if there exists a \(\kk\)-affinoid algebra \(A_V\) and a morphism of \(\kk\)-affinoid algebras \(\theta: A \to A_V\) satisfying the following universal property:
        for every bounded homomorphism of \(\kk\)-affinoid algebras \(\psi: A \to B\) such that the induced map on spectra \(\scrM(\psi): \scrM(B) \to X\) has its image contained in \(V\), 
        there exists a unique bounded homomorphism \(\varphi: A_V \to B\) such that the following right diagram commutes:
        \[
        \begin{tikzcd}
            & V  \arrow[dl, hook'] & \\
            X  & & \scrM(B) \arrow[ul] \arrow[ll, "\scrM(\psi)"]
        \end{tikzcd} 
        \quad \leadsto \quad
        \begin{tikzcd}
            & A_V  \arrow[dr, gray, "\varphi"] & \\
            A \arrow[ur, "\theta"] \arrow[rr, "\psi"'] & & B
        \end{tikzcd}
        \]
        In this case, we say that \(V\) is represented by the \(\kk\)-affinoid algebra \(A_V\).
        If \(A_V\) is strict, then we say that \(V\) is a strict affinoid domain.
    \end{definition}
    \begin{slogan}
        A closed subset \(V \subset X\) is an affinoid domain if the functor ``\(\Mor(-, V)\)'' is representable by a \(\kk\)-affinoid algebra \(A_V\).
    \end{slogan}

    % \Yang{Why we consider closed subset rather that open subset?}

    There are several important classes of affinoid domains.

    \begin{construction}\label{constr:weierstrass_Laurent_rational_domains}
        Let \(f=(f_1,\ldots,f_n),g = (g_1,\ldots,g_m)\) be two tuples of elements in \(A\) and \(r=(r_1,\ldots,r_n),s=(s_1,\ldots,s_m)\) be two tuples of positive real numbers.
        Consider the following closed subset of \(X\):
        \[ X\left(\underline{f/r};\underline{s/g}\right) \coloneqq \left\{ x \in X \colon\ |f_i(x)| \le r_i, |g_j(x)| \ge s_j, 1 \le i \le n, 1 \le j \le m \right\}. \]
        Such a closed subset is called a \emph{Laurent domain} of \(X\).
        If \(m=0\), then it is called a \emph{Weierstrass domain} of \(X\), denoted by 
        \[ X\left(\underline{f/r}\right) \coloneqq \left\{ x \in X \colon |f_i(x)| \le r_i, 1 \le i \le n \right\}. \]

        More generally, let \(f=(f_1,\ldots,f_n),g\) be elements in \(A\) such that the ideal generated by them is the whole algebra \(A\).
        Set \(p=(p_1,\ldots,p_n)\) be a tuple of positive real numbers.
        We define the following closed subset of \(X\):
        \[ X\left(\underline{(f/g)/p}\right) \coloneqq \left\{ x \in X \colon |f_i(x)| \le p_i |g(x)|, 1 \le i \le n \right\}. \]
        Such a closed subset is called a \emph{rational domain} of \(X\).
    \end{construction}  

    \begin{proposition}\label{prop:weierstrass_Laurent_rational_domains_are_affinoid_domains}
        Let \(A\) be a \(\kk\)-affinoid algebra, and let \(X = \scrM(A)\) be the associated affinoid space.
        Then any Weierstrass domain, Laurent domain, or rational domain of \(X\) is an affinoid domain.
        Moreover, if \(A\) is strict, then these affinoid domains are all strict.
        \Yang{To be checked.}
    \end{proposition}
    \begin{proof}
        \Yang{To be added.}
    \end{proof}

    We have a sequence of inclusion:
    % \[ \{\text{Weierstrass domains}\} \subseteq \{\text{Laurent domains}\} \subseteq \{\text{Rational domains}\} \subseteq \{\text{Affinoid domains}\}. \]
    \[ \{\text{Affinoid domains}\} \supseteq \{\text{Rational domains}\} \supseteq \{\text{Laurent domains}\} \supseteq \{\text{Weierstrass domains}\}. \]

    \Yang{Why Laurent domain is rational domain?}

    \Yang{Is these domain open?}
    
    \begin{proposition}\label{prop:spectrum_of_affinoid_algebra_represent_affinoid_domain}
        Let \(V \subseteq X\) be an affinoid domain represented by the \(\kk\)-affinoid algebra \(A_V\).
        Then the natural map on spectra \(\scrM(\theta): \scrM(A_V) \to X\) induced by the representing homomorphism \(\theta: A \to A_V\) has its image equal to \(V\), and induces a homeomorphism \(\scrM(A_V) \cong V\).
    \end{proposition}
    \begin{proof}
        \Yang{To be added.}
    \end{proof}

    \begin{proposition}\label{prop:affinoid_domain_is_flat_over_base}
        Let \(V \subseteq X\) be an affinoid domain represented by the \(\kk\)-affinoid algebra \(A_V\).
        Then the natural bounded homomorphism \(\varphi: A \to A_V\) is flat.
    \end{proposition}
    \begin{proof}
        \Yang{To be added.}
    \end{proof}

    \begin{proposition}\label{prop:inverse_image_of_affinoid_domain}
        Let \(f: Y=\scrM(B) \to X = \scrM(A)\) be a morphism of \(\kk\)-affinoid spaces induced by a bounded homomorphism of \(\kk\)-affinoid algebras \(\phi: A \to B\).
        Let \(V \subseteq X\) be an affinoid domain represented by the \(\kk\)-affinoid algebra \(A_V\).
        Then the inverse image \(f^{-1}(V) \subseteq Y\) is an affinoid domain represented by the \(\kk\)-affinoid algebra \(B \hat{\otimes}_A A_V\).
        \Yang{To be revised.}
    \end{proposition}
    \begin{proof}
        \Yang{To be added.}
    \end{proof}

    \begin{proposition}\label{prop:intersection_of_affinoid_domain}
        Let \(V_1,V_2 \subseteq X\) be two affinoid domains represented by the \(\kk\)-affinoid algebras \(A_{V_1}, A_{V_2}\), respectively.
        Then the intersection \(V_1 \cap V_2 \subseteq X\) is an affinoid domain represented by the \(\kk\)-affinoid algebra \(A_{V_1} \hat{\otimes}_A A_{V_2}\).
        \Yang{To be checked.}
    \end{proposition}
    \begin{proof}
        \Yang{To be added.}
    \end{proof}


    \begin{proposition}\label{prop:strictly_affinoid_domain_form_a_system_of_neighborhood}
        Let \(x \in X\) be a point.
        Then the collection of strict affinoid domains of \(X\) containing \(x\) forms a fundamental system of neighborhoods of \(x\).
        Explicitly, for any open neighborhood \(U \subseteq X\) of \(x\), there exists a strict affinoid domain \(V \subseteq X\) such that \(x \in V \subseteq U\).
        \Yang{To be checked.}
    \end{proposition}




    \begin{theorem}[Gerritzen-Grauert theorem]\label{thm:gerritzen_grauert_theorem}
        Let \(V \subseteq X\) be an affinoid domain.
        Then there exist finitely many rational domains \(V_1,\ldots,V_m \subseteq X\) such that \(V = V_1 \cup \cdots \cup V_m\).
        \Yang{To be checked.}
    \end{theorem}
    \begin{proof}
        \Yang{To be added.}
    \end{proof}


    
    \begin{definition}\label{def:dimension_of_affinoid_domain}
        Let \(V \subseteq X\) be an affinoid domain represented by the \(\kk\)-affinoid algebra \(A_V\).
        Suppose that \(A_V \cten \KK\) is strictly \(\KK\)-affinoid for some completed valued field extension \(\KK/\kk\).
        The \emph{dimension} of \(V\) is defined as the Krull dimension of the affinoid algebra \(A_V \cten \KK\):
        \[ \dim(V) := \dim_{\text{Krull}}(A_V \cten \KK). \]
        % \Yang{To be revised.}
    \end{definition}

    By \cref{prop:krull_dimension_of_affinoid_domain_under_ground_field_changing}, the definition is independent of the choice of \(\KK\).



\subsection{The Grothendieck topology and structure sheaf}

    \begin{theorem}[Acyclicity theorem]\label{thm:acyclicity_theorem_for_affinoid_domains}
        Let \(A\) be a \(\kk\)-affinoid algebra, and let \(X = \scrM(A)\) be the associated affinoid space.
        Let \(\{V_i\}_{i \in I}\) be a finite covering of \(X\) by affinoid domains.
        For each finite subset \(J \subseteq I\), set \(V_J := \bigcap_{j \in J} V_j\), which is an affinoid domain by \cref{prop:intersection_of_affinoid_domain}.
        Then the augmented Čech complex
        \[
        0 \to A \to \prod_{i \in I} A_{V_i} \to \prod_{\substack{J \subseteq I \\ |J| = 2}} A_{V_J} \to \prod_{\substack{J \subseteq I \\ |J| = 3}} A_{V_J} \to \cdots
        \]
        is exact.
        \Yang{To be checked.}
    \end{theorem}

    \begin{definition}\label{def:special_domain}
        A \emph{special domain} of \(X\) is a finite union of affinoid domains of \(X\).
        Let \(V \subseteq X\) be a special domain, say \(V = V_1 \cup \cdots \cup V_m\) with each \(V_i\) an affinoid domain.
        We define 
        \[ A_V \coloneqq \ker\left( \prod_{i=1}^m A_{V_i} \to \prod_{1 \le i < j \le m} A_{V_i \cap V_j} \right). \]
        \Yang{To be checked.}
    \end{definition}

    \begin{proposition}\label{prop:when_is_special_domain_affinoid_domain}
        Let \(V \subseteq X\) be a special domain, say \(V = V_1 \cup \cdots \cup V_m\) with each \(V_i\) an affinoid domain.
        Then \(V\) is an affinoid domain if and only if the \(\kk\)-algebra \(A_V\) defined in \cref{def:special_domain} is a \(\kk\)-affinoid algebra.
        In this case, \(V\) is represented by the \(\kk\)-affinoid algebra \(A_V\).
        \Yang{To be added.}
    \end{proposition}

    \begin{definition}\label{def:site_of_special_domains_on_affinoid_space}
        The \emph{site of special domains} on \(X\), is defined as follows:
        \begin{itemize}
            \item the underlying category has as objects the special domains \(V \subseteq X\), and the morphisms are given by the inclusion maps;
            \item a covering of an objects \(V \subseteq X\) is a finite family of affinoid domains \(\{V_i\}_{i \in I}\) such that \(V = \bigcup_{i \in I} V_i\).
        \end{itemize}
        The \emph{structure sheaf on site of special domains} is given by assigning to each special domain \(V \subseteq X\) the ring \(A_V\).
    \Yang{To be revised.}
    \end{definition}

    \begin{definition}\label{def:structure_sheaf_on_affinoid_space}
        For every open subset \(U \subseteq X\), we define the ring of analytic functions on \(U\) as follows:
        \[ \scrO_X(U) \coloneqq \varprojlim_{V \subseteq U} A_V \ \in \Alg_{\kk}, \]
        where the limit is taken over all special domains \(V \subseteq U\), and \(A_V\) is the \(\kk\)-algebra representing the special domain \(V\).
    \Yang{To be revised.}
    \end{definition}


    \begin{remark}\label{rmk:two_different_site_and_structure_sheaf_for_affinoid_spaces}
        Note that in \cref{def:site_of_special_domains_on_affinoid_space} we have defined a site on \(X\) consisting of special domains, and in \cref{def:structure_sheaf_on_affinoid_space} we have defined a structure sheaf on the usual topological space \(X\).
        In fact, these two definitions are closely related.
        \Yang{To be revised.}
    \end{remark}

    \begin{definition}\label{def:the_category_of_k_affinoid_space}
        The \emph{category of \(\kk\)-affinoid spaces} is defined as the opposite category of the category of \(\kk\)-affinoid algebras:
        \[ {\kk} \coloneqq (\Alg_{\kk}^{\text{aff}})^{\op}. \]
        \Yang{To be revised.}
    \end{definition}

    \Yang{Coherent sheaf.}








\subsection{Boundary and interior}

    Let \(A, B\) be two \(\kk\)-affinoid algebras, and let \(f: Y=\scrM(B) \to X = \scrM(A)\) be a morphism of \(\kk\)-affinoid spaces induced by a bounded homomorphism of \(\kk\)-affinoid algebras \(\phi: A \to B\).

    \begin{definition}\label{def:relative_boundary_and_interior}
        The \emph{relative interior} of \(Y\) over \(X\) is defined as
        \[ (Y/X) := \{ y \in Y \colon \text{there exists an admissible surjective} A \to B \text{ inducing } y \}. \]
        \Yang{To be revised.}
    \end{definition}

    \begin{definition}\label{def:Shilov_boundary}
        
    \end{definition}


\subsection{Some local properties}

    \Yang{The local ring \(\scrO_{X,x}\).}

    \Yang{The residue field and its completion.}

    \begin{example}\label{eg:residue_field_of_affinoid_space_depends_on_neighborhood}
        \Yang{To be added.}
    \end{example}

    \begin{proposition}\label{prop:local_ring_of_affinoid_space_is_noetherian}
        The local ring \(\scrO_{X,x}\) at any point \(x \in X\) is noetherian.
    \end{proposition}


    \begin{definition}\label{def:local_properities_of_affinoid_space}
        We say that the affinoid space \(X\) has a local property \(\mathcal{P}\) at a point \(x \in X\) if the local ring \(\scrO_{X,x}\) has the property \(\mathcal{P}\).
        \Yang{To be added.}
    \end{definition}